% !TeX root = ../../main.tex
\subsection{関数のグラフと変域}

関数の問題では，次の3つの視点を意識すると立式しやすい.

\begin{techniquebox}[tech:01]{\textbf{グラフを読む3つのポイント}}
    \begin{enumerate}
        \item \textbf{傾き} … 右上がりか，右下がりか（$a$ の符号）
        \item \textbf{切片・通過点} … グラフがどの点を通るか（代入して確認）
        \item \textbf{変域} … $x$ の範囲が決まったとき，$y$ がとりうる範囲を考える
    \end{enumerate}
\end{techniquebox}

\begin{techniquebox}[tech:02]{\textbf{比例定数を求める}}
    正比例・反比例では，与えられた1組の値を代入して比例定数 $a$ を求める.
    \begin{itemize}
        \item 正比例：$y = ax$ に代入
        \item 反比例：$y = \dfrac{a}{x}$ または $xy = a$ に代入
    \end{itemize}
\end{techniquebox}

\begin{techniquebox}[tech:03]{\textbf{一次関数の式を求める}}
    傾き $a$ と1点がわかれば $y = ax + b$ が定まる.
    2点がわかれば，2つの方程式を連立して $a$，$b$ を求める.
\end{techniquebox}

\begin{example}[【例題 1】]
    一次関数のグラフが点 $(1,\,4)$，$(3,\,10)$ を通る. この一次関数の式を求めなさい.
\end{example}

\begin{proof}\mbox{}\\
    $y = ax + b$ とおく.

    \begin{hiddenanswer}
        2点を代入すると，
        \begin{align*}
            4 &= a + b \\
            10 &= 3a + b
        \end{align*}
        連立して解くと $a = 3$，$b = 1$.
        \finalanswer{答え：$y = 3x + 1$}
    \end{hiddenanswer}
\end{proof}

\newpage
